\documentclass{article}
\usepackage[T1]{fontenc}
\usepackage{lmodern}
\usepackage{graphicx}
\usepackage{amsmath,amssymb}
\usepackage[a4paper,margin=2cm]{geometry}
\usepackage[absolute,overlay]{textpos}
\setlength{\TPHorizModule}{1mm}
\setlength{\TPVertModule}{1mm}
\usepackage[indonesian]{babel}
\usepackage{hyperref}
\setlength{\emergencystretch}{2em}
% unit: Halaman 54

\begin{document}

% === Halaman 54 (python/native) ===

54

Menurut ad-Durayhim, jenis-jenis cipher dapat dikelompokkan ke dalam delapan 

tipe:

(1) transposisi, 

(2) substitusi,

(3) penambahan atau reduksi jumlah huruf, 

(4) penggunaan piranti sandi,

(5) penggantian huruf dengan angka yang diboboti secara desimal, 

(6) penyandian huruf dengan menggunakan kata-kata, 

(7) penggantian huruf dengan nama generik, 

(8) menggunakan simbol atau tanda untuk menyatakan huruf.

Cryptology was born among Arabs. They were the first to 

discover and write down the methods of cryptanalysis.

(David Kahn – Penulis buku: The Code Breaker)

\end{document}
